\begin{table}
\caption{Original IV and ML-IV (DML-PLIV) Estimates}
\label{tab:iv_vs_mliv_table}
\begin{tabular}{lcc}
\toprule
Estimator & Column 3 & Column 4 \\
\midrule
Original & -0.286* (0.153) & -0.221*** (0.078) \\
Linear & -0.255*** (0.070) & -0.200*** (0.055) \\
Ridge & -0.250*** (0.074) & -0.211*** (0.075) \\
LASSO & -0.272*** (0.085) & -0.186** (0.082) \\
Elastic Net & -0.281*** (0.088) & -0.199*** (0.077) \\
Polynomial Ridge & -0.160 (0.138) & -0.122 (0.092) \\
SVR & -0.218** (0.105) & -0.119* (0.070) \\
Kernel Ridge & -0.283** (0.113) & -0.181** (0.074) \\
Random Forest & -0.243* (0.141) & -0.114 (0.116) \\
Extra Trees & -0.223*** (0.083) & -0.148 (0.095) \\
Gradient Boosting & -0.286** (0.143) & -0.170 (0.107) \\
\bottomrule
\\[-0.5em]
\multicolumn{3}{p{0.95\linewidth}}{\footnotesize Notes: Entries are coefficients on log slave exports per unit land area. Original rows reproduce Nunn (2008) with the conventional standard errors printed in the paper. ML rows use 5-fold cross-fitting repeated 500 times, with one common fold partition shared by all nuisance functions within a repeat. Coefficients are medians across repeats; standard errors use the median aggregation of Chernozhukov et al. (2018), $\widehat{se} = \sqrt{\mathrm{med}_r[se_r^2 + (\hat\theta_r - \tilde\theta)^2]}$, combining within-split uncertainty and split variation. * p<0.10, ** p<0.05, *** p<0.01.}\\
\end{tabular}
\end{table}
